\documentclass[11pt]{article}

% --- Packages ---
\usepackage[top=1.2in, bottom=1.2in, left=1.25in, right=1.25in]{geometry}
\usepackage[hidelinks]{hyperref}
\usepackage{booktabs}
\usepackage[table]{xcolor}
\usepackage{microtype}
\usepackage{parskip}
\usepackage{enumitem}
\usepackage{array}
\usepackage{fancyhdr}

% --- Page style ---
\pagestyle{fancy}
\fancyhf{}
\fancyhead[L]{\small\itshape Sim-to-Stage: A Paradigm for Live Performance Production}
\fancyhead[R]{\small Position Paper v0.1}
\fancyfoot[C]{\thepage}
\renewcommand{\headrulewidth}{0.4pt}

% --- List spacing ---
\setlist{itemsep=3pt, topsep=4pt, parsep=0pt}

% --- Colors ---
\definecolor{highlight}{rgb}{1.0, 0.97, 0.80}

% --- Title ---
\title{%
  \textbf{Sim-to-Stage:}\\[4pt]
  A Paradigm for Live Performance Production%
}

\author{%
  Tiantian Chen\\[2pt]
  \small Sim-to-Stage Research\\
  \small \href{https://simtostage.com}{\texttt{simtostage.com}}
}

\date{%
  Position Paper v0.1 $\cdot$ May 14, 2026\\[4pt]
  \small License: Creative Commons Attribution 4.0 International (CC BY 4.0)\\[2pt]
  \small Repository: \href{https://github.com/sim-to-stage/sim-to-stage}%
    {\texttt{github.com/sim-to-stage/sim-to-stage}}
}

\begin{document}

\maketitle
\thispagestyle{plain}

% -----------------------------------------------------------------------
\begin{abstract}
We introduce \textbf{Sim-to-Stage}: a paradigm for compiling director
intent into rehearsable, executable, audit-traceable live-stage artifacts
under a single authoritative truth source.  Sim-to-Stage is to live
performance production what \textbf{Sim-to-Real} is to
robotics---a discipline for training, validating, and transferring complex
multi-modal artifacts from a simulation environment to a real-world target.
In Sim-to-Stage, the real-world target is the live stage and the live
audience.  We define the \textbf{Sim-to-Stage gap} as the paradigm's
central technical challenge and decompose it into seven open sub-gaps.
We propose seven \textbf{architecture primitives}---one useful decomposition
that a mature Sim-to-Stage system is likely to instantiate variants of---
abstracted from an early reference implementation in development at StageR.
This position paper is a starting coordinate, not a solved problem; we
explicitly invite external research groups, theater companies, and
live-production industrial implementers to extend, contest, and improve
the paradigm.
\end{abstract}

\tableofcontents
\newpage

% -----------------------------------------------------------------------
\section{Introduction}

A live theatrical performance is a coordinated multi-modal artifact:
script, blocking, lighting cues, sound cues, video projections, scenic
transitions, performer interpretation, audience interaction---all unfolding
in real time over an evening of performance.  Producing such an artifact at
high quality has historically been organized through human craft alone:
directors, designers, technicians, and performers iterating across weeks of
rehearsal, accumulating implicit knowledge that lives in production journals,
cue sheets, and individual memory.

In the past three years, three forces have arrived simultaneously that
change the underlying engineering question.

First, large language models and multi-modal generative models have become
sufficiently capable to participate as collaborators in theatrical creative
work---not merely as tools for narrow tasks, but as departmental specialists.

Second, the robotics community has formalized \textbf{Sim-to-Real} as a
paradigm for transferring policies trained in simulation to physical
deployment, complete with mitigation techniques (domain randomization,
system identification, residual learning) and a vocabulary of ``gap''
diagnosis.

Third, the AI infrastructure community has converged on \textbf{world
models}---large-scale simulations of physical or social environments---as
a primary research direction.

These three forces converge on a question: \textbf{what would it look like
to produce a live performance the way modern robotics produces a deployed
policy?}  That is---to simulate the production end-to-end in a
director-gated environment, to validate the simulated rehearsal against a
single authoritative truth source, and to compile the validated simulation
into one-shot executable artifacts that human performers and technical crew
deploy on the live stage.  We call this paradigm \textbf{Sim-to-Stage}.

This paper makes four contributions:

\begin{enumerate}
\item We define the Sim-to-Stage paradigm precisely
      (Section~\ref{sec:paradigm}).
\item We construct a structural \textbf{homomorphism} between Sim-to-Real
      (robotics) and Sim-to-Stage (live performance), showing the two are
      not merely analogous in metaphor but structurally homologous as
      engineering pipelines (Section~\ref{sec:homomorphism}).
\item We define the \textbf{Sim-to-Stage gap} as the paradigm's central
      open problem and decompose it into seven distinct sub-gaps that admit
      independent research (Section~\ref{sec:gap}).
\item We propose seven \textbf{architecture primitives}---one useful
      decomposition that a mature Sim-to-Stage system is likely to
      instantiate variants of (Section~\ref{sec:primitives}).
\end{enumerate}

We are deliberate about what this paper does \textbf{not} claim.  We do
not claim the Sim-to-Stage gap has been closed.  We do not claim the
reference implementation has solved the seven primitives in their general
form.  We do not claim the paradigm subsumes all live-performance
production---only that it provides a useful frame for one mode of
production-augmented work.  Section~\ref{sec:discussion} details
limitations and open questions in full.

% -----------------------------------------------------------------------
\section{Background}

\subsection{The Sim-to-Real Paradigm}

The Sim-to-Real paradigm originates in robotics, where the cost and risk of
training control policies directly on physical robots motivated a
transfer-learning approach: train policies in a physics simulator (Isaac
Sim, MuJoCo, Habitat, CARLA), then deploy to physical hardware.  The
central technical problem is the \textbf{Sim-to-Real gap}---the discrepancy
between the simulator and the physical world that causes a policy trained in
sim to fail on real hardware.

Three families of mitigation are well-established.  \textbf{Domain
randomization} trains the policy across a distribution of simulator
parameters (mass, friction, lighting) so the real world appears as one of
many seen environments.  \textbf{System identification} estimates the
parameters of the physical world to bring the simulator into closer
correspondence with reality.  \textbf{Residual policy learning} trains a
corrective policy on top of the sim-trained policy that learns to account
for residual sim-to-real discrepancy.

The Sim-to-Real paradigm has become canonical in robotics because it makes
a previously intractable engineering question tractable: it gives the field
a \emph{vocabulary} (the gap, the mitigations), a \emph{benchmark
structure} (sim eval + transfer eval), and an \emph{architectural pattern}
(sim trainer, transfer step, deployment policy).  This paper proposes that
live performance production stands today in a similar position to robotics
circa 2015, and that an analogous paradigm---Sim-to-Stage---provides
similar vocabulary, structure, and pattern.

\subsection{Live Performance Production Today}

Many productions unfold over weeks or months of rehearsal, technical
preparation, and previews.  The artifact that ultimately runs on stage is
the joint product of dozens of specialists: director, dramaturg, lighting
designer, sound designer, video designer, scenic designer, costume
designer, stage manager, technical director, and a cast of performers.
The coordination problem alone is substantial; the creative problem is, by
traditional account, the entirety of theatrical art.

Several technological currents already exist in live-production practice.
Stage management software (QLab, ETC Eos, Disguise) captures cue scripts
and operates them at performance time.  Digital previz tools (Capture,
WYSIWYG, depence\textsuperscript{2}) allow lighting and projection
designers to preview their work before installation.  Collaborative document
tools and shared cue sheets manage cross-department coordination.  However,
these tools operate as \textbf{departmental islands}: each specialist's
tool optimizes that specialist's work in isolation, with handoffs across
department boundaries handled through unstructured human communication.

The Sim-to-Stage paradigm proposes a different architecture: a \textbf{single
director-gated truth source} representing the production, against which all
specialist work is reconciled; a \textbf{multi-modal simulation
environment} in which the production can be rehearsed end-to-end across
modalities; and a \textbf{compiler} that translates the validated simulation
into one-shot executable artifacts deployable by the human performers and
technical crew.

\subsection{Why ``Sim-to-Stage'' Now}

Three preconditions for Sim-to-Stage have arrived approximately together
in 2024--2026:

\begin{enumerate}
\item \textbf{Multi-modal generative models} are now capable of producing
      draft artifacts (lighting plots, sound compositions, video sequences,
      scenic visualizations) at quality sufficient for director review and
      iteration.  This makes specialist agents tractable.
\item \textbf{Long-context and persistent-memory language models} can hold
      an entire production's worth of script, notes, blocking, and rationale
      in working memory, enabling a Director Intelligence Profile that
      remembers across rehearsals.
\item \textbf{Sim-to-Real} as a paradigm has matured to the point that its
      vocabulary and architectural lessons (gap diagnosis, mitigation
      families, benchmark structure) can be ported to adjacent transfer
      domains.
\end{enumerate}

The Sim-to-Stage paradigm therefore arrives at a specific historical moment,
and we believe the moment is genuinely open: we are not aware of a widely
adopted canonical paradigm for this domain.

% -----------------------------------------------------------------------
\section{The Sim-to-Stage Paradigm}
\label{sec:paradigm}

\subsection{Definition}

\textbf{Sim-to-Stage} is a paradigm for producing live performances in
which:

\begin{enumerate}
\item The production is represented as an \textbf{authoritative,
      director-gated, multi-modal artifact} (the Authoritative Truth
      Source---Section~\ref{sec:ats}) maintained as a single source of
      truth across all departments.
\item The production is \textbf{rehearsed end-to-end in a simulation
      environment} that models, to varying fidelity, lighting state, sound
      state, video state, scenic state, blocking, and performance timing.
\item Specialist work is contributed by \textbf{department-level agents}
      (human or AI) whose proposals are reconciled against the truth source
      through \textbf{director-mediated confirmation gates}.
\item The validated simulation is \textbf{compiled into one-shot executable
      artifacts} (cue scripts, packages, briefings) that human performers
      and technical crew deploy on the live stage.
\item Every decision, every gate, every revision is recorded in a
      \textbf{provenance ledger} that maintains audit traceability from
      director intent to live execution.
\end{enumerate}

A Sim-to-Stage system is one that instantiates variants of these five
properties.  An early reference implementation is StageR
(\url{https://simtostage.com}), but the paradigm is intended to be larger
than any single implementation, and we invite external groups to
instantiate it independently.

\subsection{Homomorphism with Sim-to-Real}
\label{sec:homomorphism}

Sim-to-Stage is not merely analogous to Sim-to-Real in metaphor; the two
paradigms exhibit a structural \textbf{homomorphism}---a systematic
correspondence of roles between the two transfer pipelines.
Table~\ref{tab:homomorphism} documents the correspondence.

\begin{table}[htbp]
\centering
\caption{The Sim-to-Real $\leftrightarrow$ Sim-to-Stage Homomorphism}
\label{tab:homomorphism}
\small
\renewcommand{\arraystretch}{1.3}
\begin{tabular}{p{2.9cm} p{3.6cm} p{6.2cm}}
\toprule
\textbf{Aspect} &
\textbf{Sim-to-Real (Robotics)} &
\textbf{Sim-to-Stage (Live Performance)} \\
\midrule
Target domain        & Physical world                   & Live stage / live audience \\
Simulator            & Physics engine (Isaac Sim, MuJoCo) & Stage simulation environment \\
Object of transfer   & Robot policy / control program   & Cue script + production package \\
Source of truth      & Reward function                  & Director intent (codified, gated) \\
Validation           & Simulated rollout                & Simulated rehearsal \\
\rowcolor{highlight}
Transfer challenge   & Sim-to-real gap                  & \textbf{Sim-to-stage gap (this paper)} \\
Mitigation: random.  & Domain randomization             & Variability rehearsal \\
Mitigation: ident.   & System identification            & Performer / venue identification \\
Mitigation: residual & Residual policy learning         & Dress-rehearsal residual adjustment \\
Executor             & Robot embodiment                 & Human performers + technical crew \\
Re-run semantics     & Deploy, observe, redeploy        & One-shot per show \\
Success criterion    & Task-completion metric           & Director judgment + audience response \\
Multi-agent struct.  & Single robot or fleet            & Irreducibly multi-department (lighting,
                                                         sound, video, scenic, blocking, script)
                                                         plus cast \\
Governance           & Safety review, deployment gate   & Director-gated Confirmation Gates,
                                                         technical and dress rehearsal sign-offs \\
\bottomrule
\end{tabular}
\end{table}

We claim this homomorphism is more than a stylistic flourish: it identifies
an engineering correspondence we believe is load-bearing.  Each row of
Table~\ref{tab:homomorphism} maps a Sim-to-Real role to a Sim-to-Stage role
that occupies the same structural position in the transfer pipeline.  We
argue the homomorphism licenses the import of Sim-to-Real's accumulated
engineering wisdom---randomization-style mitigation, identification-style
mitigation, residual-style mitigation, deployment-gate governance---into
Sim-to-Stage practice.  The remaining sections develop the specifics.

% -----------------------------------------------------------------------
\section{The Sim-to-Stage Gap}
\label{sec:gap}

\subsection{Definition}

We define the \textbf{Sim-to-Stage gap} as the family of discrepancies
between a simulated rehearsal and the live performance that cause a
Sim-validated production package to fail on the live stage.  Like the
Sim-to-Real gap in robotics, the Sim-to-Stage gap is not a single
phenomenon but a structural collection of distinct sub-gaps, each amenable
to its own analysis and mitigation.  Closing the Sim-to-Stage gap is the
central open research program of the paradigm.

\subsection{Seven Identified Sub-Gaps}

We identify seven sub-gaps.  We do not claim the list is exhaustive; we
claim it is \emph{useful}: each sub-gap is independently studyable, and
progress on any one is progress on the paradigm.

\paragraph{Sub-gap 1 $\cdot$ The Timing-Variability gap.}
Simulated rehearsal runs against a synthetic clock; live performance runs
against human performance timing, which varies with audience response,
performer state, and unrepeatable contingency.  A cue script that fires
precisely on a synthetic clock will fail when human performance is faster,
slower, or restructured by ad-lib.  \textit{Candidate direction:} robust
cue compilation that admits a tolerance window per cue, with sim-time
domain randomization across performance-timing distributions.

\paragraph{Sub-gap 2 $\cdot$ The Director-Authority gap.}
In Sim-to-Real, the reward function is the authority---a formal
mathematical object.  In Sim-to-Stage, the director's intent is the
authority, but intent is informal: aesthetic, stylistic, intuitive,
evolving through rehearsal.  How does a system formalize an informal
authority sufficiently that the rest of the pipeline can act on it, without
flattening the aesthetic content?  \textit{Candidate direction:} the
Director Intelligence Profile (Section~\ref{sec:dip})---a model of
director preference that combines explicit declared rules, learned
stylistic patterns, and decision-by-decision confirmation.

\paragraph{Sub-gap 3 $\cdot$ The Multi-Modal Coherence gap.}
A live performance is a joint artifact across lighting, sound, video,
scenic, blocking, script, and performer interpretation.  Each modality has
its own conventions, file formats, and specialist tools.  Simulating one
modality well is tractable; simulating all coherently is harder, and
preserving coherence under transfer is harder still.
\textit{Candidate direction:} specialist production agents
(Section~\ref{sec:specialists}) coordinating against a shared Authoritative
Truth Source (Section~\ref{sec:ats}), with cross-modal coherence checks at
each Confirmation Gate.

\paragraph{Sub-gap 4 $\cdot$ The Human-Execution gap.}
Robots execute the deployed policy autonomously.  In Sim-to-Stage, the
deployed artifact is executed by humans---actors interpret blocking, stage
managers call cues, technicians operate consoles.  Each human introduces
interpretation, slippage, and modification.  How does a sim-validated
artifact survive contact with human execution?
\textit{Candidate direction:} compilation targets that explicitly model
human execution affordances (cue scripts written for human operators;
rehearsable, not just executable), and previews / dress rehearsals as
transfer-validation checkpoints.

\paragraph{Sub-gap 5 $\cdot$ The One-Shot Deployment gap.}
A robot policy can fail in deployment, be observed, and be redeployed
after a fix.  A live performance is one-shot per show: there is no retry.
Sim-validation must be sufficient for one-shot deployment, a regime closer
to safety-critical Sim-to-Real (surgical robotics, aerospace) than to
typical robotics.  \textit{Candidate direction:} layered defense---sim
eval, then specialist eval, then technical rehearsal eval, then
dress-rehearsal eval, then preview eval---with each layer admitting only
what the previous layer validated.

\paragraph{Sub-gap 6 $\cdot$ The Long-Horizon Narrative gap.}
A robot policy may have a horizon of seconds to minutes; a live performance
typically has a horizon of an hour or more of continuous interconnected
causal narrative (foreshadowing, character arcs, dramatic structure,
payoff).  Local Sim-validation of each scene does not imply global
validation of long-horizon coherence.  \textit{Candidate direction:}
narrative-aware Project Intelligence (a long-context understanding
component) coupled with explicit dramaturgical evaluation criteria during
sim rehearsal.

\paragraph{Sub-gap 7 $\cdot$ The Aesthetic-Judgment gap.}
Robot success criteria are predominantly functional (the robot completed
the task or not).  Live performance success criteria are partially
aesthetic---was the moment moving, did the comedy land, did the audience
lean in.  Even formal metrics (run time, sound level, cue precision)
underdetermine aesthetic success.  \textit{Candidate direction:} the
Director Intelligence Profile as aesthetic judge
(Section~\ref{sec:dip}), combined with structured audience response
signals in previews; explicit acknowledgment that aesthetic judgment will
remain partially human throughout the paradigm's lifetime.

These seven sub-gaps are the proposed map of the Sim-to-Stage research
territory.  We invite extension, refinement, and challenge.

% -----------------------------------------------------------------------
\section{Architecture Primitives}
\label{sec:primitives}

We propose the following seven architecture primitives as one useful
decomposition of a Sim-to-Stage system.  We expect a mature Sim-to-Stage
system to instantiate variants of these primitives, but we do not claim
this decomposition is unique or exhaustive
(see Section~\ref{sec:discussion}).  We define each primitive at the
paradigm level---the role it plays, the interface it presents to the rest
of the system---without committing to implementation details, which are
properly the subject of later technical reports.

\subsection{The Authoritative Truth Source (ATS)}
\label{sec:ats}

The ATS is a single, canonical, append-only registry of the production
state: the script as currently authorized, the blocking as currently
authorized, the cue list as currently authorized, the design choices as
currently authorized.  Every other component of the system reads from the
ATS and proposes against it; only Confirmation Gates
(Section~\ref{sec:gate}) admit writes to it.  The ATS is the formal
codification of ``the production as the director currently authorizes it.''

\begin{quote}
\textit{Role in the homomorphism:} The ATS is the Sim-to-Stage analog of
the policy parameter set in Sim-to-Real---the canonical object being
trained, validated, and transferred.
\end{quote}

\subsection{The Director Intelligence Profile (DIP)}
\label{sec:dip}

The DIP is the system's model of the director: their declared aesthetic
rules, their stylistic patterns inferred from prior decisions, their
preferences for ambiguity resolution, and their authority delegation policy
(what they decide themselves vs.\ what specialists may decide).  The DIP
is the mechanism by which the informal authority of director intent is
rendered partially formal---enough to drive specialist proposals and Sim
evaluation, without flattening aesthetic judgment.

\begin{quote}
\textit{Role in the homomorphism:} The DIP is the Sim-to-Stage analog of
the reward function in Sim-to-Real.  Where robotics often gets a single
mathematical reward function, Sim-to-Stage gets a partial, evolving,
gated stand-in.
\end{quote}

\subsection{Specialist Production Agents}
\label{sec:specialists}

Specialist agents are domain-specific generative and orchestration
components: a lighting agent that drafts lighting plots and cues, a sound
agent for sound, a video agent for video, a scenic agent for scenic, a
blocking agent for blocking, and so on.  Each specialist agent reads the
current ATS, proposes changes against it, and submits proposals to the
Confirmation Gate for the relevant department.  Specialists may be AI
agents, human professionals using AI-augmented tools, or any combination.

\begin{quote}
\textit{Role in the homomorphism:} Specialist agents are the Sim-to-Stage
analog of the multi-agent or multi-skill policy structures emerging in
robotics fleet learning.  Live performance is irreducibly multi-department,
so this multi-agent structure is structurally necessary, not stylistic.
\end{quote}

\subsection{The Confirmation Gate}
\label{sec:gate}

A Confirmation Gate is the director-mediated decision point at which
specialist proposals are evaluated and either accepted into the ATS,
rejected, or returned for revision.  Confirmation Gates may be high-touch
(the director reviews each proposal individually) or low-touch (specialists
operate within director-declared autonomy bands and only flag edge cases),
but they are always present as the authorization mechanism.  The
Confirmation Gate is the formal site at which director authority enters
the pipeline.

\begin{quote}
\textit{Role in the homomorphism:} The Confirmation Gate is the
Sim-to-Stage analog of the deployment gate in safety-critical Sim-to-Real
practice---the explicit moment at which ``this artifact is authorized for
the next stage.''
\end{quote}

\subsection{The Stage Simulation Environment}
\label{sec:sim-env}

The Stage Simulation Environment is the sim itself: a multi-modal model of
the live stage at some fidelity---lighting state, sound state, video state,
scenic state, blocking, performer timing.  The Sim Environment runs the
validated ATS as a simulated rehearsal, surfacing problems (cross-modal
incoherence, dramatic dead spots, technical conflicts) before they appear
in physical rehearsal.

\begin{quote}
\textit{Role in the homomorphism:} The Sim Environment is the Sim-to-Stage
analog of the physics simulator in Sim-to-Real.  It is the substrate of
the paradigm; the rest of the system is defined relative to it.
\end{quote}

\subsection{The Cue Compiler}
\label{sec:compiler}

The Cue Compiler is the component that translates the Sim-validated ATS
into one-shot executable artifacts for live execution: cue scripts callable
by stage management, console programs deployable to lighting / sound /
video consoles, briefing packages for cast and technical crew.  The Cue
Compiler is responsible for materializing the sim's validated state into
formats that live execution can act on.

\begin{quote}
\textit{Role in the homomorphism:} The Cue Compiler is the Sim-to-Stage
analog of the policy export step in Sim-to-Real---the moment at which the
sim-internal representation is converted into a deployable artifact.
\end{quote}

\subsection{The Provenance Ledger}
\label{sec:provenance}

The Provenance Ledger is an append-only audit trail of every decision in
the production: every specialist proposal, every Confirmation Gate outcome,
every ATS revision, every Sim rehearsal verdict.  The ledger enables
accountability (``who decided this and why''), reproducibility (``what was
the state of the ATS at this point''), and downstream learning (training
the DIP on the director's actual revealed decisions over time).

\begin{quote}
\textit{Role in the homomorphism:} The Provenance Ledger is the
Sim-to-Stage analog of the experiment log and model registry in modern
robotics MLOps practice.
\end{quote}

% -----------------------------------------------------------------------
\section{Toward a Sim-to-Stage Benchmark}

A paradigm is established not only by its definition but by the benchmark
community against which progress is measured.  We sketch the structure of a
Sim-to-Stage benchmark here; the full proposal will land separately as
\texttt{benchmark.md} in the canonical repository.

We propose that a Sim-to-Stage benchmark should evaluate, at minimum:

\begin{itemize}
\item \textbf{Cross-modal coherence}---does the Sim-rehearsed production
      exhibit consistent lighting / sound / video / scenic / blocking
      coordination at each cue point?
\item \textbf{Transfer fidelity}---does the executed live performance match
      the Sim-rehearsed state on observable dimensions (cue timing,
      cross-modal coordination, executed-vs-planned divergence)?
\item \textbf{Director-intent retention}---does the executed production
      preserve the director's declared aesthetic and stylistic priorities
      as recorded in the DIP?
\item \textbf{Audit traceability}---for any property of the executed
      production, can the Provenance Ledger trace it back to a declared
      director decision?
\item \textbf{Production efficiency}---does the paradigm reduce the
      wall-clock effort to producible state, relative to the unaided
      baseline?
\end{itemize}

A serious Sim-to-Stage benchmark will require a benchmark \emph{corpus}
of productions across genres, scales, and traditions.  We invite
participation in defining this corpus.

% -----------------------------------------------------------------------
\section{Related Work}

The Sim-to-Stage paradigm draws on, and is meant to extend, several
adjacent literatures.

\begin{itemize}
\item \textbf{Sim-to-Real in robotics} is the immediate ancestor; the
      homomorphism in Section~\ref{sec:homomorphism} is the formal
      statement of the inheritance.
\item \textbf{World models} in machine learning provide simulation substrate
      techniques applicable to Sim-to-Stage rehearsal environments.
\item \textbf{Multi-agent systems} literature informs the specialist
      production agent layer.
\item \textbf{Theater production scholarship}---production journals, design
      process literature, the history of stage management---provides the
      ground truth against which any Sim-to-Stage system must be evaluated.
\item \textbf{Live production tooling}---QLab, ETC Eos, Disguise, Capture,
      WYSIWYG, depence\textsuperscript{2}---represents the
      departmental-island state of the practice that Sim-to-Stage proposes
      to unify under a shared truth source.
\end{itemize}

We will expand related-work coverage in subsequent revisions.

% -----------------------------------------------------------------------
\section{Discussion and Limitations}
\label{sec:discussion}

We are explicit about what this position paper does \textbf{not} claim and
where the paradigm's open problems remain.

\begin{itemize}
\item \textbf{The seven sub-gaps are unsolved.}  We have named them, not
      closed them.  Each is a research direction, not a result.
\item \textbf{The seven architecture primitives are not unique.}  Other
      decompositions of a Sim-to-Stage system are possible.  We claim ours
      is \emph{useful}, not \emph{exclusive}.
\item \textbf{Aesthetic judgment will remain partially human.}  Sub-gap~7
      is, in our view, partially irreducible; we do not propose the paradigm
      fully formalizes theatrical taste, only that the DIP and the
      Confirmation Gate provide the right interfaces for human aesthetic
      judgment to enter the pipeline.
\item \textbf{The paradigm has not been validated at scale.}  The reference
      implementation StageR is in development.  We are honest that this
      paper is, at its publication date, ahead of empirical evidence.
      Subsequent papers in this repository will report on empirical progress
      as it accumulates.
\item \textbf{The paradigm may not cover all live production.}  Some
      traditions---devised work, improvisation-heavy forms,
      audience-participatory performance---may fit awkwardly into the ATS /
      Confirmation Gate / Cue Compiler structure.  We do not claim universal
      applicability.
\end{itemize}

% -----------------------------------------------------------------------
\section{Conclusion}

We have introduced the Sim-to-Stage paradigm: a structural homomorphism of
Sim-to-Real adapted to the domain of live performance production.  We have
defined the Sim-to-Stage gap as the paradigm's central open problem and
decomposed it into seven sub-gaps.  We have proposed seven architecture
primitives as one useful decomposition that a mature Sim-to-Stage system
will likely instantiate variants of.  We have sketched the structure of a
Sim-to-Stage benchmark.

We are at the beginning of this research program, not the end.  We invite
extensions, refinements, challenges, and alternative formulations from
external research groups, theater companies, and live-production industrial
implementers.  The paradigm belongs to its community; this paper is a
starting coordinate.

% -----------------------------------------------------------------------
\section*{How to Cite This Work}
\addcontentsline{toc}{section}{How to Cite This Work}

If you reference the term ``Sim-to-Stage'' or this paradigm in academic
or industrial work, please cite:

\begin{verbatim}
@misc{chen2026simtostage,
  title  = {Sim-to-Stage: A Paradigm for Live Performance Production},
  author = {Chen, Tiantian},
  year   = {2026},
  url    = {https://simtostage.com},
  note   = {Position paper v0.1.
            Canonical repository:
            https://github.com/sim-to-stage/sim-to-stage}
}
\end{verbatim}

A machine-readable \texttt{CITATION.cff} is provided in the canonical
repository and powers GitHub's ``Cite this repository'' auto-export.

\section*{License}
\addcontentsline{toc}{section}{License}

This work is licensed under
\href{https://creativecommons.org/licenses/by/4.0/}{Creative Commons
Attribution 4.0 International (CC BY 4.0)}.  You may use, share, adapt,
translate, and build upon this material---including for commercial
purposes---provided you give appropriate credit.  The paradigm belongs to
its community; the license invites contribution.

\section*{Versioning}
\addcontentsline{toc}{section}{Versioning}

This is Position Paper \textbf{v0.1}.  We expect substantial revision in
v0.2 and v0.3 as the paradigm matures, additional empirical evidence
accumulates, and external contributors challenge or extend the formulation.
We welcome issues and pull requests on the canonical repository:
\url{https://github.com/sim-to-stage/sim-to-stage}.

\end{document}
